\documentclass[11pt]{scrartcl}
\usepackage[polish]{mystd}
\title{Deskryptory sieciowe}
\subtitle{jako narzędzia rozróżniające ludzkie i zwierzęce sieci społecznościowe}
\author{Michał Dobranowski}
\date{czerwiec 2023}

\usepackage[backend=biber,style=ieee]{biblatex}
\addbibresource{refs.bib}

\begin{document}
\maketitle
\tableofcontents
\eject

\section{Wstęp i oznaczenia}
Jeśli nie zaznaczono inaczej, to przez $V$ będziemy oznaczać zbiór $n$ wierzchołków w grafie, a przez $E$ zbiór $m$ krawędzi. Ponadto, $\deg(v)$ oznacza stopień wierzchołka $v$ (czyli liczbę jego sąsiadów).

\subsection{Deskryptory sieciowe}
\paragraph*{Rozkład stopni wierzchołków}
\[ P(k) = \frac{n_k}{n}, \]
gdzie $n_k$ to wierzchołków o stopniu $k$.

\paragraph*{Clustering coefficient}
\[ C(v) = \frac{T(v)}{\deg(v)(\deg(v) - 1)}, \]
gdzie $T(v)$ to liczba trójkątów $K_3$ zawierających $v$. Jest miarą tego, jak wiele sąsiadów danego wierzchołka jest połączonych; pozwala ocenić, czy w sieci występują grupy wierzchołków połączonych ze sobą (klastry).

\paragraph*{Closeness centrality}
\[ C_C(v) = \frac{1}{\sum\limits_{t \in V} d_G(v, t)}, \]
gdzie $d_G(v, t)$ to odległość od wierzchołka $v$ do $t$ (czyli długość najkrószej ścieżki między nimi). Określa, jak blisko znajduje się dany węzeł do wszystkich innych węzłów w sieci, a więc jak łatwo z danego wierzchołka możemy osiągnąć inne wierzchołki.

\paragraph*{Betweenness centrality}
\[ C_B(v) = \sum_{s\neq v\neq t} \frac{\sigma_{st}(v)}{\sigma_{st}}, \]
gdzie $\sigma_{st}$ to liczba najkrótszych ścieżek między wierzchołkami $s, t$, a $\sigma_{st}(v)$ to liczba najkrótszych ścieżek między wierzchołkami $s, t$, które przechodzą przez wierzchołek $v \neq s, t$. Ta miara opisuje jak istotny jest dany węzeł w komunikacji między innymi węzłami w sieci. W porównaniu do poprzednich deskryptorów, jego optymalne wyliczenie nie jest tak proste, stosuje się do tego celu algorytm Brandesa \cite{Brandes2001}.

\section{Deskryptory w modelach sieci losowych}
W tej sekcji przedstawimy trzy standardowe modele losowych sieci oraz porównamy opisane deskryptory na ich przykładzie. Sieci zostały wygenerowane za pomocą biblioteki NetworkX \cite{NetworkX} języka Python. Wszystkie wykresy powstały z użyciem matplotlib \cite{matplotlib} oraz seaborn \cite{seaborn}.

\subsection{Model Erdősa-Rényi'ego}
Najprostrzym modelem losowych grafów (oraz najstarszym), jest ten rozważany przez Paula Erdősa oraz Alfréda Rényi'ego \cite{ErdosRenyi59}. Każda para wierzchołków jest jest połączona krawędzią z jednakowym prawdopodobieństem $p$.

\subsection{Model Wattsa-Strogatza}
Model przedstawiony przez Duncana Wattsa i Stevena Strogatza \cite{WattsStrogatz98} charakteryzuje się małą odległością między wierzchołkami (proporcjonalną do $\log n$) oraz wysoką klasteryzacją (ma cechy sieci typu \textit{small-world}). Algorytm tworzenia takiej sieci wygląda następująco: zaczynamy od grafu stanowiącego $\frac{k}{2}$-tą potęgę cyklu $C_n$. Następnie przepinamy każdą krawędź do innego wierzchołka (wybranego losowo) z prawdopodobieństwem $\beta$.


\subsection{Model Barabásiego-Albert}
Zaproponowany przez Alberta-László Barabásiego i Rékę Albert \cite{BarabasiAlbert99} algorytm generowania losowych grafów tworzy sieci bezskalowe (\textit{scale-free}) to znaczy takie, w których rozkład stopni wierzchołków jest rozkładem potęgowym. Tego typu sieci tworzy się przez inkrementowanie liczby krawędzi, w każdym kroku faworyzując wierzchołki o dużym stopniu (to znaczy, że prawdopodobieństwo przyłączenia nowej krawędzi do danego wierzchołka jest wprost proporcjonalne do jego stopnia).
W ten sposób otrzymujemy
\[ P(k) \sim k^{-3}. \]




Network Repository \cite{NetworkRepository}. Animal Social Network Repository \cite{AnimalSocialNetworkRepository}.

\printbibliography[heading=bibintoc]
\end{document}